\input{preamble.tex}

\title{Отчет по лабораторной работе № 13. \\ Статическая маршрутизация в Интернете. Планирование}
\author{Сущенко Алина Николаевна \\ НПИбд-01-23}
\date{2026}

\begin{document}

\maketitle
\newpage

\tableofcontents

\newpage
\section{Цель работы}
Провести подготовку по организации взаимодействия через сеть провайдера посредством статической маршрутизации локальной сети с сетью основного здания, расположенного в 42-м квартале в Москве, и сетью филиала, расположенного в г. Сочи. (Рис.\ref{01}) 
\newpage

\section{Выполнение работы}
\subsection{Изменение схемы сети}
\begin{enumerate}
\item Внесем изменения в нашу схему в соответствии с L1. Мы добавили квартал 42 в Москве и г. Сочи. Присволили везде названия в соответствии с соглашением об именовании.
\begin{center}
    \centering
    \includegraphics[width=\textwidth]{../images/01.png}
    \captionof{figure}{Обновленная схема сети}
    \label{01}
\end{center}

\item Далее, на медиаконвертерах  замените имеющиеся модули на \texttt{PT-REPEATER-NM-1FFE} и \texttt{PT-REPEATER-NM-1CFE} для подключения витой пары по технологии \texttt{Fast Ethernet} и оптоволокна соответственно. (Рис.\ref{02}):
\begin{center}
    \centering
    \includegraphics[width=\textwidth]{../images/02.png}
    \captionof{figure}{Замена модулей на медиаконвертерах}
    \label{02}
\end{center}

\subsection{Первоначальная настройка маршрутизатора ansusenko-q42-gw-1}
\item Произведем начальную настройку маршрутизатора \texttt{ansusenko-q42-gw-1} (Рис.\ref{03}):
\begin{center}
    \centering
    \includegraphics[width=\textwidth]{../images/03.png}
    \captionof{figure}{Настройка маршрутизатора.}
    \label{03}
\end{center}

\item На маршрутизатора заранее добавили дополнительный интерфейс NM-2FE2W.(Рис.\ref{04}):
\begin{center}
    \centering
    \includegraphics[width=\textwidth]{../images/04.png}
    \captionof{figure}{Добавленный дополнительный интерфейс.}
    \label{04}
\end{center}

\item Затем продолжили настройку как того требует задание.(Рис.\ref{09}):

\begin{center}
    \centering
    \includegraphics[width=\textwidth]{../images/09.png}
    \captionof{figure}{Окончание настройки маршрутизатора.}
    \label{09}
\end{center}

\subsection{Добавление городов и зданий в физической схеме проекта}
\item Для начала в городе Москва добавили здание отвечающее за квартал-42.(Рис.\ref{05}):

\begin{center}
    \centering
    \includegraphics[width=\textwidth]{../images/05.png}
    \captionof{figure}{Схема города Москва.}
    \label{05}
\end{center}

\item В здании 42-го квартала разместили наши объекты.(Рис.\ref{06}):

\begin{center}
    \centering
    \includegraphics[width=\textwidth]{../images/06.png}
    \captionof{figure}{Размещение объектов в основном здании 42-го квартала в Москве.}
    \label{06}
\end{center}

\item Так же проверим здание \texttt{Provider} на наличие новых медиаконвертеров.(Рис.\ref{07}):

\begin{center}
    \centering
    \includegraphics[width=\textwidth]{../images/07.png}
    \captionof{figure}{Размещение объектов в здании Provider.}
    \label{07}
\end{center}

\item Так же создали город \texttt{Sochi} и разместили в нем объекты.(Рис.\ref{08}):

\begin{center}
    \centering
    \includegraphics[width=\textwidth]{../images/08.png}
    \captionof{figure}{Размещение объектов в здании филиала в г. Сочи.}
    \label{08}
\end{center}

\subsection{Первоначальная настройка коммутатора ansusenko-q42-sw-1}
\item Для коммутатора \texttt{ansusenko-q42-sw-1} провели первоначальную настройку.(Рис.\ref{10}):

\begin{center}
    \centering
    \includegraphics[width=\textwidth]{../images/10.png}
    \captionof{figure}{Начальная настройка коммутатора ansusenko-q42-sw-1.}
    \label{10}
\end{center}

\subsection{Первоначальная настройка маршрутизирующего коммутатора ansusenko-hostel-gw-1}
\item Для маршрутизирующего коммутатора \texttt{ansusenko-hostel-gw-1} провели первоначальную настройку.(Рис.\ref{11}):

\begin{center}
    \centering
    \includegraphics[width=\textwidth]{../images/11.png}
    \captionof{figure}{Начальная настройка ansusenko-hostel-gw-1.}
    \label{11}
\end{center}

\subsection{Первоначальная настройка коммутатора ansusenko-hostel-sw-1}
\item Для коммутатора \texttt{ansusenko-hostel-sw-1} провели первоначальную настройку.(Рис.\ref{12}):

\begin{center}
    \centering
    \includegraphics[width=\textwidth]{../images/12.png}
    \captionof{figure}{Начальная настройка ansusenko-hostel-sw-1.}
    \label{12}
\end{center}

\subsection{Схема подключения подсети филиала в г. Сочи}
\item Для коммутатора \texttt{sch-ansusenko-sw-1} провели первоначальную настройку.(Рис.\ref{13}):

\begin{center}
    \centering
    \includegraphics[width=\textwidth]{../images/13.png}
    \captionof{figure}{Начальная настройка sch-ansusenko-sw-1.}
    \label{13}
\end{center}

\item Для маршрутизатора \texttt{sch-ansusenko-gw-1} провели первоначальную настройку.(Рис.\ref{14}):

\begin{center}
    \centering
    \includegraphics[width=\textwidth]{../images/14.png}
    \captionof{figure}{Начальная настройка sch-ansusenko-gw-1.}
    \label{14}
\end{center}

\end{enumerate}

\section{Выводы}
В результате лабораторной работы мы усложнили нашу сеть, успешно провели подготовку по организации взаимодействия через сеть провайдера посредством статической маршрутизации локальной сети с сетью основного здания, расположенного в 42-м квартале в Москве, и сетью филиала, расположенного в г. Сочи.

\end{document}
